
\vspace{-2mm}
\section{Preliminaries}
\label{sec:preliminaries}
Consider a Temporal Knowledge Graph (TKG) $\mathcal{G} = (\mathcal{E}, \mathcal{R}, \mathcal{T})$, where $\mathcal{E}$, $\mathcal{R}$, and $\mathcal{T}$ represent the set of entities, relations, and timestamps, respectively. 
$\mathcal{G} = (\mathcal{E}, \mathcal{R}, \mathcal{T})$ contains abundant temporal factual knowledge in the form of quadruples, i.e., $\mathcal{G} = \{(e_h, r, e_t, t) \mid e_h,e_t \in \mathcal{E}, r \in \mathcal{R}, t \in \mathcal{T} \}$.
% where $s$ and $o$ denote subject and object entities, $p$ is the predicate, and $t$ is the temporal validity of the fact. 
% 
% Each fact captures not only structural semantics but also temporal specificity.
Temporal constraint defines a condition related to a specific time point or interval that must be satisfied by both the answer and its supporting evidence. Following Allen's interval algebra~\cite{allen1984towards}, this includes the 13 classical temporal relations, as well as extended operators such as temporal set relations, duration comparisons, and sorting mechanisms (e.g., first, last, nth)~\cite{TimelineKGQA}. Temporal constraints thus enable fine-grained alignment between natural language queries and temporally consistent evidence paths.




\begin{definition}[Temporal Path] 
Given a TKG $\mathcal{G}$, a temporal  path is a connected sequence of temporal facts, represented as:  
$
path_{\mathcal{G}}(e_1, e_{l+1}) = \{(e_1, r_1, e_2, t_1), (e_2, r_2, e_3, t_2), \ldots, (e_l, r_l, e_{l+1}, t_l)\},$ where $l$ denotes the length of the path, i.e., ${\rm{length}}({\rm{path}}_{\mathcal{G}}(e_1, e_{l+1})) = l$. The path must satisfy two constraints:  
(1) connectivity: the tail entity of each fact is identical to the head entity of the next;  
(2) temporal monotonicity: timestamps are non-decreasing, \ie $t_1 \leq t_2 \leq \cdots \leq t_l$.  

% Thus, a temporal reasoning path guarantees that the structural sequence and chronological order of evidence are both respected.
\end{definition}


\begin{example}[Temporal Path]
Consider a temporal path between ``Merkel'' and ``EU'' with a length of 3 :
$path_{\mathcal{G}}(\text{Merkel}, \text{EU}) = \{
(\text{Merkel},~\textit{visit},~\text{Paris},~2012),~
(\text{Paris},~\textit{host},~\text{Conference},~2013),$ $
(\text{Conference},~\textit{attended\_by},~\text{EU},~2014)
\}$,
and can be visualized as:
\[
\text{Merkel} \xrightarrow[\text{2012}]{\textit{visit}} 
\text{Paris} \xrightarrow[\text{2013}]{\textit{host}} 
\text{Conference} \xrightarrow[\text{2014}]{\textit{attended\_by}} 
\text{EU}.
\]
% \noindent It is a valid temporal path because entities are directly connected and timestamps are monotonic (2012 $\le$ 2013 $\le$ 2014). 
% The length of the path is 3.
\end{example}

\begin{definition}[Temporal Reasoning Path]  
Given a TKG $\mathcal{G}$, and an entity list $[e_1,e_2,\ldots,e_n]$, a temporal reasoning path is a sequence of temporal segments
$
TRP_{\mathcal{G}}([e_1,\ldots,e_n]) = \{P_1,P_2,\ldots,P_{n}\},
$
where each $P_i = path_{\mathcal{G}}(e_i,n_i)$ for some $n_i \in \mathcal{E}$,  
and the global monotonicity condition holds:
$
t^{end}(P_i) \leq t^{start}(P_{i+1}), \quad \forall~1 \leq i < n.
$
\end{definition}



\begin{example}[Temporal Reasoning Path]
Consider a temporal reasoning path $TRP_{\mathcal{G}}([{Obama}, {Beijing}, {Paris}])$, by two temporally aligned segments:
% $TRP_{\mathcal{G}}([{Obama}, {Beijing}, {Paris}])$ 
% = {$Path_{\mathcal{G}}$(Obama, UN), $Path_{\mathcal{G}}$(Beijing, EU) } 
$P_1$= \{(Obama, meet, UN, 2009)\}, $P_2$= \{(Beijing, linked\_via, EU, 2011), (EU, event\_in, Paris, 2012)\}, visualized as:
\[
\text{Obama} \xrightarrow[\text{2009}]{\textit{meet}} 
\text{UN} 
\;\;\leadsto\;\;
\text{Beijing} \xrightarrow[\text{2011}]{\textit{linked\_via}} 
\text{EU} \xrightarrow[\text{2012}]{\textit{event\_in}} 
\text{Paris}.
\]
% \noindent It is a valid temporal reasoning path because the segments are temporally ordered 
% ($2009 \le 2011 \le 2012$), even though they are not directly edge-contiguous in $\mathcal{G}$.
\end{example}


% \begin{definition}[Temporal Reasoning Path]
% Given a temporal KG $\mathcal{G}$ and an entity pairs list $list$ = $\{(s,t)|\forall s, t \in \mathcal{E}\}$, a temporal reasoning path is defined as a concatenation of temporal paths between consecutive entities:
% $
% path_{\mathcal{G}}(list) = \{ path_{\mathcal{G}}(e_1,e_2)\}.
% $
% Unlike a temporal reasoning path, the start and end entities of $list_e$ need not be directly connected. The essential requirement is that all involved facts maintain monotonic temporal order across segments, \ie if the last fact of $path^{temp}_{\mathcal{G}}(e_i, e_{i+1})$ has time $t^{(i)}_{end}$, then $t^{(i)}_{end} \leq t^{(i+1)}_{start}$ for the subsequent segment. This definition enables reasoning across multiple disjoint entities under strict chronological consistency.
% \end{definition}

% \begin{definition}[Temporal Reasoning Path]
% Given a TKG $\mathcal{G}$ and an entity list 
% $list_e = [e_1, e_2, \ldots, e_l]$.
% A temporal reasoning path is a sequence of temporal segments that connects 
% consecutive entities in $list_e$, wich denoted as $
% Path{\mathcal{G}}(list_e) = \{ (e_i \leadsto e_{i+1}) \mid 1 \leq i < l \},
% $
% where each link $e_i \leadsto e_{i+1}$ is realized by a pair of temporally aligned paths 
% via some intermediate anchors.  
% % through possibly intermediate anchors. 
% Formally, for each consecutive pair $(e_i, e_{i+1})$, there exist intermediate entities 
% $u, v \in \mathcal{E}$, i.e.,
% $
% path_{\mathcal{G}}(e_i, u) \quad \text{and} \quad path_{\mathcal{G}}(v, e_{i+1})
% $
% and 
% $
% t^{end}(path_{\mathcal{G}}(e_i, u)) \;\leq\; t^{start}(path_{\mathcal{G}}(v, e_{i+1})).
% $
% \end{definition}


% \begin{definition}[Temporal Reasoning Path]
% Given a TKG $\mathcal{G}$ and a list of entities
% $list_e = [e_1, e_2, \ldots, e_l]$, 
% a temporal reasoning path is an ordered sequence of temporal links between consecutive entities, denoted as  
% $TRP_{\mathcal{G}}(list_e) = \{ e_i \leadsto e_{i+1} \mid 1 \leq i < l \}$.
% Each link $e_i \leadsto e_{i+1}$ may be realized through a intermediate entity, e.g., $path_{\mathcal{G}}(e_i,u)$ or $path_{\mathcal{G}}(u,e_{i+1})$ with $u \in \mathcal{E}$, such that their timestamps satisfy  
% $t^{end}(path_{\mathcal{G}}(e_i,u)) \leq t^{start}(path_{\mathcal{G}}(v,e_{i+1}))$.
% \end{definition}

% \begin{definition}[Temporal Reasoning Path]
% Given a TKG $\mathcal{G}$ and an ordered entity list 
% $list_e = [e_1, e_2, \ldots, e_l]$, 
% a temporal reasoning path (TRP) is an ordered sequence of temporal links between consecutive entities, denoted as  
% $TRP_{\mathcal{G}}(list_e) = \{ e_i \leadsto e_{i+1} \mid 1 \leq i < l \}$.

% Each link $e_i \leadsto e_{i+1}$ is not required to be a direct edge. Instead, it is realized through a single intermediate anchor $u \in \mathcal{E}$, such that both $path_{\mathcal{G}}(e_i,u)$ and $path_{\mathcal{G}}(u,e_{i+1})$ exist, and their timestamps satisfy  
% $t^{end}(path_{\mathcal{G}}(e_i,u)) \leq t^{start}(path_{\mathcal{G}}(u,e_{i+1}))$.  

% Equivalently, a TRP can be represented as a list of entity pairs 
% $[(e_i,u), (u,e_{i+1})]$ for each $(e_i,e_{i+1})$, where local temporal monotonicity is preserved and the global chain remains chronologically consistent.  
% \end{definition}
% \begin{definition}[Temporal Reasoning Path]
% Let $\mathcal{G}=(\mathcal{E},\mathcal{R},\mathcal{T},\mathcal{F})$ be a temporal knowledge graph.  
% A temporal reasoning path (TRP) is an ordered list of entity pairs
% \[
% TRP = [(s_1,t_1), (s_2,t_2), \ldots, (s_k,t_k)], \quad s_i,t_i \in \mathcal{E}.
% \]

% Each pair $(s_i,t_i)$ is associated with a temporal path 
% $path_{\mathcal{G}}(s_i,t_i) = \{(v_0,r_1,v_1,\tau_1),\ldots,(v_{m-1},r_m,v_m,\tau_m)\}$ 
% with non-decreasing timestamps.  

% The global requirement is temporal monotonicity: if $\tau^{end}(i)$ and $\tau^{start}(i+1)$ are the end and start times of consecutive paths, then
% \[
% \tau^{end}(i) \leq \tau^{start}(i+1), \quad \forall~1 \leq i < k.
% \]

% Thus, a TRP is a sequence of temporally ordered reasoning segments, 
% which need not be structurally adjacent in $\mathcal{G}$, but must align chronologically.
% \end{definition}

% \begin{definition}[Temporal Reasoning Path]
% Given a TKG $\mathcal{G}$ and an ordered entity list 
% $list_e = [e_1, e_2, \ldots, e_n]$, 
% a temporal reasoning path (TRP) is defined as a set of temporally ordered path segments
% \[
% TRP_{\mathcal{G}}(list_e) = \{ path_{\mathcal{G}}(e_i,n_i) \mid 1 \leq i \leq n \},
% \]
% where each $n_i \in \mathcal{E}$ is an intermediate entity (possibly $n_i = e_{i+1}$).  

% The temporal monotonicity condition requires that
% \[
% t^{end}(path_{\mathcal{G}}(e_i,n_i)) \;\leq\; t^{start}(path_{\mathcal{G}}(e_{i+1},n_{i+1})), 
% \quad \forall~ 1 \leq i < n.
% \]

% Thus, a TRP connects the given entity list not necessarily by structural adjacency, 
% but by aligning their associated path segments in chronological order.
% \end{definition}



% \begin{definition}[Temporal Reasoning Path]
% Given a temporal KG $\mathcal{G}$ and a list of entity pairs 
% $list = [(s_1,t_1), (s_2,t_2), \ldots, (s_k,t_k)]$ with $s_i,t_i \in \mathcal{E}$, 
% a temporal reasoning path is defined as an ordered collection of temporal paths between each pair:
% \[
% TRP_{\mathcal{G}}(list) = \{ path_{\mathcal{G}}(s_1,t_1),\; path_{\mathcal{G}}(s_2,t_2),\; \ldots,\; path_{\mathcal{G}}(s_k,t_k) \}.
% \]

% Each $path_{\mathcal{G}}(s_i,t_i)$ is a temporal path satisfying connectivity and local temporal monotonicity.  
% The global requirement is that these segments respect chronological order across pairs, \ie 
% if the last fact in $path_{\mathcal{G}}(s_i,t_i)$ has timestamp $t^{(i)}_{end}$, 
% and the first fact in $path_{\mathcal{G}}(s_{i+1},t_{i+1})$ has timestamp $t^{(i+1)}_{start}$, 
% then
% \[
% t^{(i)}_{end} \leq t^{(i+1)}_{start}, \quad \forall i < k.
% \]

% Unlike a single temporal path, a temporal reasoning path may connect disjoint entity pairs, 
% but it must preserve strict chronological consistency across all segments.
% \end{definition}

\noindent
Temporal Knowledge Graph Question Answering (TKGQA)
is a fundamental reasoning task based on TKGs. Given a natural language question $Q$ and a TKG $\mathcal{G}$, the objective is to devise a function $f$ that identify answer entities or timestamps $a \in Answer(Q)$ utilizing knowledge in $\mathcal{G}$, \ie $a = f(q, \mathcal{G})$.
% Following previous works \cite{sun2019pullnet, rogluo2023reasoning, tog1.0sun2023think, tog2.0ma2024think}, 
Consistent with previous research \cite{sun2019pullnet, rogluo2023reasoning, tog1.0sun2023think, tog2.0ma2024think},
we assume topic entities $Topics$ are mentioned in $Q$ and answer entities or timestamps  $Answer(q)$ in ground truth are linked to $\mathcal{G}$, i.e., $Topics \subseteq \mathcal{E} \text{ and }  Answer(Q) \subseteq \{\mathcal{E} \cup \mathcal{T}\}$.
% \end{problem}

% \section{Preliminaries}
% \label{sec:preliminary}

% \paragraph{Temporal Knowledge Graph.}
% A Temporal Knowledge Graph (TKG) is denoted as $\mathcal{G} = (\mathcal{E}, \mathcal{R}, \mathcal{T}, \mathcal{F})$, 
% where $\mathcal{E}$ is the set of entities, $\mathcal{R}$ is the set of relations, and $\mathcal{T}$ is the set of temporal markers (e.g., timestamps or intervals). 
% The fact set $\mathcal{F}$ consists of quadruples $(s, r, o, t)$ where $s, o \in \mathcal{E}$, $r \in \mathcal{R}$, and $t \in \mathcal{T}$, each representing that subject $s$ and object $o$ are linked by relation $r$ at time $t$. 
% Compared with a static KG, the temporal dimension in $\mathcal{G}$ enforces that reasoning chains must respect both structural connectivity and chronological consistency.

% \paragraph{Temporal Constraints.}
% A temporal constraint specifies a condition over a point or interval in $\mathcal{T}$ that candidate answers and supporting evidence must satisfy. 
% Typical constraints include Allen’s interval relations~\cite{allen1984} (\eg before, overlaps, during), as well as higher-level operators such as ordering (\texttt{first}/\texttt{last}), counting (\texttt{nth}), or aggregation within temporal sets~\cite{sun2025temporal}. 
% These operators ensure that inference not only identifies semantically relevant entities but also enforces time-consistent evidence paths.

% \paragraph{Temporal KGQA.}
% Given a natural language query $q$, Temporal KGQA requires finding an answer set $A(q)$ consisting of entities or timestamps by reasoning over a temporally valid subgraph of $\mathcal{G}$. 
% Formally, the task can be expressed as learning a mapping function
% \[
% f: (q, \mathcal{G}) \mapsto A(q),
% \]
% where $A(q) \subseteq \mathcal{E} \cup \mathcal{T}$. 
% Questions may include explicit temporal scopes (e.g., \textit{``Who was the president of the U.S. in 2008?''}) or implicit ones (e.g., \textit{``Who was the president during the financial crisis?''}), requiring temporal alignment between question semantics and evidence.

% \begin{definition}[Temporal Reasoning Path]
% Given a TKG $\mathcal{G}$, a temporal reasoning path is a sequence of temporal facts that connects a source entity $e_1$ to a target entity $e_{l+1}$:
% \[
% path^{temp}_{\mathcal{G}}(e_1, e_{l+1}) = \{(e_1, r_1, e_2, t_1), (e_2, r_2, e_3, t_2), \ldots, (e_l, r_l, e_{l+1}, t_l)\}.
% \]
% The path must satisfy two constraints:  
% (1) \textbf{Connectivity:} for each $i$, the object of $(e_i,r_i,e_{i+1},t_i)$ is the subject of $(e_{i+1},r_{i+1},e_{i+2},t_{i+1})$;  
% (2) \textbf{Temporal monotonicity:} timestamps are non-decreasing, i.e., $t_1 \leq t_2 \leq \cdots \leq t_l$.  
% This guarantees that reasoning evidence respects both graph structure and chronological order.
% \end{definition}

% \begin{definition}[Temporal Entity Path]
% Given $\mathcal{G}$ and a list of entities $list_e = [e_1, e_2, \ldots, e_l]$, 
% a temporal entity path is a concatenation of temporal reasoning paths between consecutive entities:
% \[
% path^{temp}_{\mathcal{G}}(list_e) = \{path^{temp}_{\mathcal{G}}(e_1, e_2), path^{temp}_{\mathcal{G}}(e_2, e_3), \ldots, path^{temp}_{\mathcal{G}}(e_{l-1}, e_l)\}.
% \]
% Unlike a single temporal reasoning path, the endpoints $e_1$ and $e_l$ are not required to be directly connected. 
% The essential requirement is that all subpaths respect temporal monotonicity, i.e., if the final timestamp of $path^{temp}_{\mathcal{G}}(e_i, e_{i+1})$ is $t^{(i)}_{end}$, then $t^{(i)}_{end} \leq t^{(i+1)}_{start}$. 
% This enables reasoning across disjoint but temporally ordered chains of evidence.
% \end{definition}



% \begin{example}[Temporal Path]
% Consider a temporal path between ``Obama'' and ``G20'' with a length of 3 :
% $path_{\mathcal{G}}(\text{Obama}, \text{G20}) = \{
% (\text{Obama},~\textit{visit},~\text{Beijing},~2009),~
% (\text{Beijing},~\textit{host},~\text{Summit},~2010),$ $
% (\text{Summit},~\textit{attended\_by},~\text{G20},~2011)
% \}$,
% and can be visualized as:
% \[
% \text{Obama} \xrightarrow[\text{2009}]{\textit{visit}} 
% \text{Beijing} \xrightarrow[\text{2010}]{\textit{host}} 
% \text{Summit} \xrightarrow[\text{2011}]{\textit{attended\_by}} 
% \text{G20}.
% \]
% % \noindent It is a valid temporal path because entities are directly connected and timestamps are monotonic (2009 $\le$ 2010 $\le$ 2011). 
% % The length of the path is 3.
% \end{example}



% \begin{example}[Temporal Reasoning Path]
% Consider a temporal reasoning path $TRP_{\mathcal{G}}([\text{Obama}, \text{Beijing}, \text{Paris}]) 
% = \{P_1, P_2\},$
% We form a TRP by aligning two temporal segments in chronological order, i.e., $TRP_{\mathcal{G}}([\text{Obama}, \text{Beijing}, \text{Paris}]) 
% = \{P_1, P_2\},$ 

% We construct the following temporal segments:
% $P_1$ = $Path_{\mathcal{G}}$
% (Obama, UN)= {(Obama , meet, UN, 2009)}, $P_2$ = $Path_{\mathcal{G}}$(Beijing, EU) = {(Beijing, linked\_via, EU, 2011),(EU, event\_in, Paris, 2012)}.


% :
% \[
% \text{Obama} 
% \xrightarrow[\text{2009}]{\textit{meet}} 
% \text{UN} 
% \;\;\leadsto\;\;
% \text{Beijing} 
% \xrightarrow[\text{2011}]{\textit{linked\_via}} 
% \text{EU} 
% \xrightarrow[\text{2012}]{\textit{event\_in}} 
% \text{Paris}.
% \]
% \noindent Here, $(\text{Obama},\textit{meet},\text{UN},2009)$ serves as a temporal bridge before the Beijing–Paris segment, 
% $\{(\text{Beijing},\textit{linked\_via},\text{EU},2011),(\text{EU},\textit{event\_in},\text{Paris},2012)\}$.
% Timestamps are monotonic ($2009 \le 2011 \le 2012$), so this is a valid TRP even though the two segments are not edge-contiguous in $\mathcal{G}$.
% \end{example}

% \begin{example}
% Consider the entity list $[e_1,e_2,e_3] = [\text{Obama}, $ $\text{Beijing}, \text{Paris}]$ in a temporal KG $\mathcal{G}$.  
% We construct the following temporal segments:
% $P_1$ = $Path_{\mathcal{G}}$
% (Obama, UN)= {(Obama , meet, UN, 2009)}, $P_2$ = $Path_{\mathcal{G}}$(Beijing, EU) = {(Beijing, linked\_via, EU, 2011),(EU, event\_in, Paris, 2012)}.
% The temporal reasoning path is
% \[
% TRP_{\mathcal{G}}([\text{Obama}, \text{Beijing}, \text{Paris}]) 
% = \{P_1, P_2\},
% \]
% with timestamps satisfying $2009 \leq 2011 \leq 2012$.  
% \end{example}